% Created 2026-04-13 lun 20:44
% Intended LaTeX compiler: pdflatex
\documentclass[11pt]{article}
\usepackage[utf8]{inputenc}
\usepackage[T1]{fontenc}
\usepackage{graphicx}
\usepackage{longtable}
\usepackage{wrapfig}
\usepackage{rotating}
\usepackage[normalem]{ulem}
\usepackage{amsmath}
\usepackage{amssymb}
\usepackage{capt-of}
\usepackage{hyperref}
\author{Abraham Nuñez}
\date{\today}
\title{Notas Desarrollo de Sistemas Web - Semestre VI}
\hypersetup{
 pdfauthor={Abraham Nuñez},
 pdftitle={Notas Desarrollo de Sistemas Web - Semestre VI},
 pdfkeywords={},
 pdfsubject={},
 pdfcreator={Emacs 30.2 (Org mode 9.7.39)}, 
 pdflang={English}}
\begin{document}

\maketitle
\tableofcontents

\section{Introducción}
\label{sec:orgfa2bf28}
\begin{itemize}
\item Profesor: Erika Meneses Rico
\item Contacto: ermeneses@uv.mx
\item Materia: Desarrollo de Sistemas Web
\item Carrera: Ciberseguridad e Infraestructura de Computo
\end{itemize}
\section{Información}
\label{sec:orga89a0cc}
\begin{itemize}
\item Tareas y practicas: 30\%
\item Proyecto: 25\%
\item Examenes: 20\%
\item Expocisión: 25\% -> 15\% revisiones y 10\% evaluación final
\end{itemize}
\section{Temario}
\label{sec:orgb9c7efa}
\subsection{Unidad 1. Fundamentos de tecnologias web}
\label{sec:orgb9a48ee}
\subsubsection{Protocolos de Internet}
\label{sec:org8b59a79}
\begin{enumerate}
\item Arquitectura TCP/IP
\label{sec:orgb26ddb8}
\begin{itemize}
\item Los equipos se resuleven mediante protocolos.
\item Los protocolos son conjuntos de normas para formatos de mensaje y procedimientos que permiten a las maquinas y los programas de aplicacion intercambiar informacion.
\item La arquitectura TCP/IP es uno de los protocolos mas viejos en los estandares de redes internas, ademas, define cuidadosamente como se meuve la informacion desde el remitente hasta el destinatario.

\begin{enumerate}
\item TCP
\begin{enumerate}
\item Es el encargado de proporcionar un servicio de comunicacion entre dos equipos con autorizacion de ambas partes y sin perdida de datos.
\item Fue creado entre 1973 a 1974 por Vint Cert y Robert Kahn.
\item Da soporte a muchas de las aplicaciones mas populares (navegadores, clientes FTP) de internet y protocolos de aplicacion (HTTP, SSH, SMTP y FTP).
\item Confiable, orientado a conexion, necesita sincronizar numeros de secuencia.
\end{enumerate}

\item IP
\begin{enumerate}
\item Es el conjunto de reglas. para enrutar y direccionar paquetes de datos para que puedan viajar a traves de las redes y llegar al destino correcto.
\end{enumerate}
\end{enumerate}

\item TCP/IP son las siglas de Transmission Control Protocol/Internet Protocol.
\item Desarrollado por DARPA
\item Usado en ARPANET en el 1971.

\begin{enumerate}
\item Diferencias entre TCP e IP
\begin{enumerate}
\item IP es la parte que obtiene la direccion a la que se envian los datos.
\item TCP se encarga de la entrega de los datos una vez hallada dicha direccion IP.
\end{enumerate}
\end{enumerate}
\end{itemize}
\begin{enumerate}
\item Funcionamiento TCP/IP
\label{sec:org5e9bb9a}
\begin{itemize}
\item Se requieren 3 cosas para la transferencia de datos:
\begin{enumerate}
\item Reglas para el empaquetamiento de datos que se envian y reciben.
\item Red interconectada.
\item Manera de encontrar la ruta para hacer llegar los datos desde la computadora que envia hasta la que recibe.
\end{enumerate}

\item TCP/IP descompone cada mensaje en paquetes que se vuelven a ensamblar en el otro extremo. Incluso, cada paquete podria tomar una ruta distinta hasta el equipo de destino si la ruta deja de estar disponible o esta muy congestionada.

\item Tambien divide tareas de comunicacion en capas. El proposito de las capas es crear un sistema estandarizado, para los distintos fabricantes de hardware y software.
\end{itemize}
\begin{enumerate}
\item UDP (User Data Protocol)
\label{sec:orgaa02bae}
\begin{itemize}
\item Es un protocolo de nivel de transporte basado en la transmision sin conexion de datagramas y representa una alternativa al protocolo TCP.
\item No establece una conexion previa.
\item Tampoco tiene confirmacion ni control de flujo, por tanto, no se sabe si ha llegado correctamente.
\item Se usa cuando se prioriza la velocidad sobre la integridad como en los protocolos DHCP, DNS o cuando se trasmite voz o video.
\end{itemize}
\item Three Hand Shake - Mensajes de TCP
\label{sec:orgb70d856}
\begin{itemize}
\item Permite verificar si el servidor esta disponible, sincronizar numeros de secuencia y establecer una sesion confiable.
\begin{itemize}
\item SYN: El cliente solicita una conexion, iniciando un numero de secuencia aleatorio.
\item SYN-ACK: El servidor responde confirmando la solicitud del cliente (SYN+1) y enviando su propio numero de secuencia inicial.
\item ACK: El cliente confirma la recepcion del SYN-ACK del servidor, incrementando el numero de secuencia del servidor.
\end{itemize}
\end{itemize}
\item Puertos
\label{sec:org4772c70}
\begin{itemize}
\item Usa puertos para permitir la comunicacion entre aplicaciones.
\item Tiene una longitud de 16 bits, por tanto, tiene un rango de 0 a 65535.
\begin{enumerate}
\item 0 - 1023: Puertos conocidos
\item 1024 - 49151: Puertos registrados
\item 49152 - 65535: Puertos privados/dinamicos
\end{enumerate}
\end{itemize}
\item Arquitectura TCP/IP
\label{sec:org45f0b5b}
\begin{itemize}
\item Application/Capa de Aplicacion: Representa datos para el usuario mas el control de codificacion y de dialogo.
\item Transport/Capa de Transporte: Admite la comunicacion entre distintos dispositivos a traves de diversas redes.
\item Internet/Capa de Internet: Determina el mejor camino a traves de una red.
\item Network interface/Capa de Acceso a la red: Controla los dispositivos del hardware y los medios que forman la red.

\item Capa y unidad
\begin{itemize}
\item Aplicacion: Datos
\item Transporte: Segmento
\item Red: Datagrama
\item Enlace: Trama
\item Fisica: Bits
\end{itemize}

\item Aplicacion: Abre una pagina web
\item Transporte: Divide el mensaje en paquetes, les pone un numero y anade verificacion de errores.
\item Internet: Anade una direccion origen y destino.
\item Red: Se envian como bits ya sea por cable o aire.
\end{itemize}
\begin{enumerate}
\item Capas del Modelo TPC-IP
\label{sec:org19f05f0}
\begin{enumerate}
\item Capa de aplicacion
\begin{itemize}
\item Define las aplicaciones de red y los servicios de internet. Estos servicios utilizan la capa de transporte para enviar y recibir datos.
\end{itemize}

\item Capa de transporte
\begin{itemize}
\item Garantiza que los paquetes lleguen en secuencia y sin errores, al intercambiar la confirmacion de la recepcion de los datos y retransmitir los paquetes perdidos (Transmision Punto a punto).
\end{itemize}

\item Capa de internet
\begin{itemize}
\item Acepta y transfiere paquetes para la red, es decir, es responsable de las funciones de direccionamiento, empaquetado y enrutamiento.
\end{itemize}
\end{enumerate}
\end{enumerate}
\end{enumerate}
\end{enumerate}
\end{enumerate}
\subsubsection{Metodologias para aplicaciones de Hipertexto}
\label{sec:org8c7d919}
\begin{enumerate}
\item Procesos para el desarrollo web
\label{sec:org3c64c1f}
\begin{itemize}
\item Se divide en dos partes:
\begin{enumerate}
\item Front-end
\item Back-end
\end{enumerate}

\item Etapas del desarrollo web:
\begin{enumerate}
\item Briefing: Documento en el que se establecen las bases de los proyectos y la informacion basica para disenar la web.
\item Planning: Se establecen los pasos para llegar a los objetos y las diferentes tareas a realizar. Se define cual es la estrategia de contenido y la estrcutura del sitio web mas adecuado para el proyecto.
\begin{itemize}
\item Existen diversas estructuras:
\begin{enumerate}
\item Lineal
\item Radial
\item Jerarquica
\item Asociativa
\end{enumerate}
\end{itemize}
\item Diseno grafico: Se hace todo el diseno (colores, tipografias, iconografia y estilos). Es importante usar wireframes (bocetos) para crear un diseno responsivo. Se pueden usar mockups (bocetos con alta fidelidad) y otros fotomontajes para mejorar la previsualizacion del diseno.
\item Determinar contenidos web: Se consideran aspectos SEO (visibilidad en buscadores) y SEM (marketing en buscadores) para el contenido, banners e informacion. Ademas de textos legales.
\item Pruebas de lanzamiento: Verifican que todo funcione correctamente ademas de ser compatible con navegadores y dispositivos.
\item Mantenimiento: Desde actualizaciones hasta nuevas funciones.
\end{enumerate}
\end{itemize}
\item Aplicaciones de hipertexto
\label{sec:orgbc23d0f}
\begin{itemize}
\item Conjunto de tecnicas y metodos organizativos que se aplcian para disenar soluciones de software informatico.
\end{itemize}
\begin{enumerate}
\item Metodologias clasicas
\label{sec:org93bf14b}
\begin{itemize}
\item Hypertext Desing Model (HDM)
\begin{itemize}
\item Una aplicacion hipermedia debe disenarse separando tres aspectos:
\begin{itemize}
\item Modelo conceptual: Tipo de informacion, entidades, relaciones, etc.
\item Modelo de navegacion: Movimiento del usuario, jerarquia de rutas.
\item Modelo de interfaz: Presentacion y organizacion de los contenidos.
\end{itemize}
\end{itemize}

\item Object-Oriented Hypermedia Design Method (OOHDM)
\begin{itemize}
\item Evolucion de HDM pero orientada a objetos con modelos independientes.
\begin{itemize}
\item Diseno conceptual: Clases y relaciones.
\item Diseno navegacional: Nodos de navegacion y enlaces.
\item Diseno abstracto de interfaz: Presentacion de los elementos.
\item Implementacion: Se implementan al sistema.
\end{itemize}
\end{itemize}

\item Web Modeling Language (WebML)
\begin{itemize}
\item Para aplicaciones web basada en modelos.
\begin{itemize}
\item Modelo de datos: Entidades y relaciones.
\item Modelo de estructura: Paginas del sitio.
\item Modelo de navegacion: Enlaces entre paginas.
\item Modelo de presentacion: Presentacion de los contenidos.
\end{itemize}
\end{itemize}
\end{itemize}
\item Metodologias modernas
\label{sec:org1c7f32b}
\begin{itemize}
\item Agiles: Desarrollo iterativo, entregas incrementales, colab con el cliente y adaptacion. Usan ciclos:
\begin{enumerate}
\item Definir
\item Disenar
\item Implementar
\item Probar
\end{enumerate}

\item Centrado al usuario (UCD/DCU): Se consideran necesidades del usuario, contexto de uso y pruebas con usuarios.
\begin{enumerate}
\item Contexto de uso
\item Especificar requerimentos
\item Diseno
\item Evaluar diseno
\end{enumerate}

\item Information Architecture: Facilitar que los usuarios encuentren y comprendan la informacion dentro de un sistema.
\begin{enumerate}
\item Sistema de organizacion
\item Sistema de navegacion
\item Sistema de etiquetado
\item Sistema de busqueda
\end{enumerate}
\end{itemize}
\end{enumerate}
\end{enumerate}
\subsection{Unidad 2. Protocolo HTTP}
\label{sec:orgd55a62c}
\subsubsection{HTTP}
\label{sec:orge148eca}
\begin{itemize}
\item Hypertext Transfer Protocol (HTTP) (o Protocolo de Transferencia de Hipertexto en español) es un protocolo de la capa de aplicación para la transmisión de documentos hipermedia, como HTML.
\item Fue diseñado para la comunicación entre los navegadores y servidores web, aunque se puede utilizar para otros propósitos también.

\item HTTP es sencillo. esta pensado y desarrollado para ser leído y fácilmente interpretado por las personas.
\item HTTP es extensible. las cabeceras de HTTP, han hecho que este protocolo sea fácil de ampliar y de experimentar con él.
\item HTTP es un protocolo con sesiones, pero sin estados. HTTP es un protocolo sin estado, es decir: no guarda ningún dato entre dos peticiones en la misma sesión. Las cookies permiten crear un contexto común para cada sesión de comunicación.
\item HTTP y conexiones. Una conexión se gestiona al nivel de la capa de trasporte, y por tanto queda fuera del alcance del protocolo HTTP.
\begin{itemize}
\item HTTP/1.0, HTTP/1.1, HTTP/2
\end{itemize}
\end{itemize}
\subsubsection{Solicitud-Respuesta}
\label{sec:org87cec2e}
\begin{enumerate}
\item El usuario final utiliza la interfaz de la aplicacióno el navegador y envía el comando al servidor a través de Internet.
\item El servidor web luego transfiere el comando al servidor solicitado.
\item El servidor solicitado busca los resultados de los comandos dados.
\item La información procesada se transfiere a la aplicación web que la envía al servidor web.
\item Ahora el servidor web proporciona los datos solicitados al usuario.
\end{enumerate}
\subsubsection{Protocolos sin estado}
\label{sec:org8714ab1}
\begin{itemize}
\item HTTP es un protocolo sin estado, esto quiere decir que el servidor no retiene información o un estatus de los usuarios durante múltiples peticiones.
\item Una sesión HTTP solo incluye una interacción petición-respuesta, la cual incluye a su vez una sesión TCP.
\item Sin embargo, algunas aplicaciones Web implementan estados o sesiones del lado del servidor utilizando tecnologías como cookies o variables ocultas dentro de formularios para almacenar información de seguimiento de usuarios.
\end{itemize}
\subsubsection{Estrcutura de las peticiones HTTP}
\label{sec:org8ea66ba}
\begin{enumerate}
\item URI
\label{sec:orgac7c388}
\end{enumerate}
\end{document}
